\documentclass[10pt,a4paper]{amsart}
\usepackage[margin=19mm]{geometry}
\usepackage{amsmath,amssymb,mathtools}
\usepackage[scr=rsfs,cal=euler]{mathalfa}
\usepackage{xfrac}
\usepackage[unicode,hidelinks,bookmarks=false]{hyperref}
\newcommand{\ie}{i.e.\ }
\newcommand{\cf}{cf.\ }
\newcommand{\resp}{respectively\ }
\newcommand{\Subsection}{\subsection}
\providecommand{\nobreakdash}{\nobreak\hbox{-}}
\setlength{\emergencystretch}{2em}
\multlinegap0pt
\usepackage{amssymb}
\usepackage[shortlabels]{enumitem}
\usepackage[normalem]{ulem}
\newtheorem{proposition}{Proposition}
\newcommand{\R}{\mathbb{R}}
\newcommand{\abs}[1]{\left\lvert #1 \right\rvert}
\newcommand{\ind}{\mathbb{1}}
\newcommand{\mF}{\mathcal{F}}
\newcommand{\mH}{\mathcal{H}}
\newcommand{\tm}{\tilde{m}}
\title{On the dimension-free control of higher order truncated Riesz transforms by higher order Riesz transforms}
\author{Maciej Kucharski, Mateusz Kwaśnicki, B{\l}a{\.z}ej Wr{\'o}bel}
\date{Source: arXiv:2507.17510v2 (CC BY 4.0)}
\begin{document}
\maketitle

\noindent\emph{Source and licence.} On the dimension-free control of higher order truncated Riesz transforms by higher order Riesz transforms. Version arXiv:2507.17510v2 (CC BY 4.0), distributed by arXiv under the Creative Commons Attribution 4.0 International licence, http://creativecommons.org/licenses/by/4.0/, as recorded in the arXiv licence metadata for this exact version and shown on the abstract page https://arxiv.org/abs/2507.17510v2. Full licence notice: \texttt{licenses/arxiv-2507.17510-CC-BY-4.0.txt}.\par\smallskip

\noindent\textit{Editorial scope.} Counted unit: the article's proposition mk (source0.tex lines 370--382). The article's complete original proof of it is delivered with it (source0.tex lines 383--447, one \texttt{proof} environment of 65 lines). No mathematics is written, repaired or re-derived: the statement and the proof are the article's own lines, in the article's own notation, and no retained line is substituted or re-flowed.  Retained uncounted as the source-local context the proof needs: lines 185--187.  Nothing was omitted from any retained span, so the package carries no omission record.  The retained lines cite 3 works, whose entries are transcribed verbatim from the article's own bibliography source in the e-print and delivered with short editorial labels recorded in the item's reference mapping.  No retained line contains a figure environment, an image-inclusion command or drawing code, so no image is delivered and the package declares no asset dependency.  Editorial formatting only, in addition: an AMS \texttt{amsart} preamble that restates the article's own declarations and macros used by the retained lines, namely the environments \texttt{proposition} on separate counters, together with the standard packages those lines need.  The retained environments print in this document as Proposition 1 in order of appearance, whereas the article prints the same items with the numbering of its own full document.  The complete native source of the article as distributed in the arXiv e-print for this exact version is bundled unchanged as \texttt{sources/182086/source0.tex}.
\begin{equation} \label{eq:R}
	R_P f(x)= \lim_{t \to 0^+} R_P^t f(x), \qquad\textrm{ where } \qquad R_P^t f(x) = \gamma_{k,d} \int_{|y|>t} \frac{P(y)}{|y|^{d+k}} f(x-y) \, dy.
\end{equation}

\begin{proposition}
\label{pro: mk}
The radial profile $m_k$ of the multiplier $\widehat{b}_k$ can be expressed by 
 \begin{equation}
\label{eq: mk 1}
m_k(r)=\frac{2^{d/2} \Gamma(\frac{d + k}{2})}{\Gamma(\frac{k}{2})} \int_{2\pi r}^\infty t^{-d/2} J_{d/2 + k - 1}(t) \, dt,\qquad r>0,
\end{equation}
and $m_k(0)=1.$ Furthermore
\begin{equation}
\label{eq: mk 2}
 m_k(r) = 1 - \frac{\Gamma(\frac{d + k}{2})}{\Gamma(\frac{d}{2} + k) \Gamma(\frac{k}{2} + 1)} \, (\pi r)^k {_1F_2}(\tfrac{k}{2}; \tfrac{d}{2} + k, \tfrac{k}{2} + 1; -(\pi r)^2),\qquad r>0.
 \end{equation}
\end{proposition}

\begin{proof}
We first justify \eqref{eq: mk 1}. Fix $P\in \mH_k$ and denote by $\rho_P^1$ the multiplier symbol of the operator $R_P^1$ given by \eqref{eq:R}. Let $\varphi(r) = \gamma_{k,d}r^{-d-k}\ind_{r > 1}$ be the radial profile of $\gamma_{k,d}|x|^{-d-k}\ind{_{|x|>1}}$. Then we have
\[
\rho_P^1(\xi)=\mF_{d}(P(\cdot)\varphi(|\cdot|))(\xi).
\]
 Using Bochner's relation for $P\in \mH_k$ (see, e.g., \cite[Corollary p.72]{stein}), together with a standard approximation argument, we see that
\[
\rho_P^1(\xi)=(-i)^k P(\xi)\Phi(\abs{\xi}),
\]
where $\Phi$ is defined by
\begin{equation}
\label{eq: defPhi}
\mF_{d+2k}(\varphi(|\cdot|))(\eta)=\Phi(|\eta|),\qquad \eta\in \R^{d+2k}.
\end{equation}
Since 
\[
\rho_P^1(\xi)=(-i)^k \frac{P(\xi)}{|\xi|^k}(|\xi|^k\Phi(\abs{\xi}))
\]
we have
\[
\mF_d[R_P^1 f](\xi)=\mF_d[R_P f](\xi)|\xi|^k\Phi(\abs{\xi}),
\]
so that
$|\xi|^k\Phi(
\abs{\xi})=\widehat{b}_k(\xi)$ and thus $m_k(r)=r^k\Phi(r),$ $r>0.$

To prove \eqref{eq: mk 1} we come back to \eqref{eq: defPhi} and write the Fourier transform of the radial function $\varphi(|y|)$ on $\R^{d+2k}$ in terms of Bessel functions (the Hankel transform). Applying \cite[Section B.5]{grafakos}, we see that
\[
 m_k(r) = r^k\frac{2\pi \gamma_{k,d}}{r^{n/2-1}} \int_1^\infty t^{-d-k} J_{n/2 - 1}(2\pi tr) t^{n/2} \, dt, 
\]
where $n=d+2k.$ Changing variables and recalling the definition of $\gamma_{k,d}$, we reach \eqref{eq: mk 1}. Since $b_k\in L^1(\R^d)$, we know that  $m_k$ is a continuous function on $[0,\infty).$ Thus the equation $m_k(0)=1$ follows from \eqref{eq: mk 1} and \href{https://dlmf.nist.gov/10.22.E43}{equation~10.22.43} in~\cite{dlmf}. 

It remains to prove \eqref{eq: mk 2}.  Let $\tm(r)=m_k(r/(2\pi))$. From \eqref{eq: mk 1} we have 
\[
    \tm'(r) = -\frac{2^{d/2} \Gamma(\frac{d + k}{2})}{\Gamma(\frac{k}{2})}\, r^{-d/2} J_{d/2 + k - 1}(r) ,
\]
and thus, by \href{https://dlmf.nist.gov/10.16.E9}{equation~10.16.9} in~\cite{dlmf},
\[
 \tm'(r) = -\frac{\Gamma(\frac{d + k}{2})}{\Gamma(\frac{d}{2} + k) \Gamma(\frac{k}{2})}  \, (\tfrac{r}{2})^{k - 1} {_0F_1}(\tfrac{d}{2} + k; -(\tfrac{r}{2})^2),
\]
where $_0F_1$ denotes the generalized hypergeometric function. Furthermore, since 
\[
 \int_0^\infty \tm'(r) \, dr = -\tm(0) =  -1 , 
\]
we have
\begin{equation*}
 \tm(r) = 1 + \int_0^r \tm'(t) \, dt = 1- \frac{\Gamma(\frac{d + k}{2})}{\Gamma(\frac{d}{2} + k) \Gamma(\frac{k}{2})} \int_{0}^r (\tfrac{t}{2})^{k - 1} {_0F_1}(\tfrac{d}{2} + k; -(\tfrac{t}{2})^2) \, dt.
\end{equation*}
Using the change of variables $u= (t/r)^2$, we obtain
\begin{equation*}
\tm(r)=1-\frac{\Gamma(\frac{d + k}{2})}{\Gamma(\frac{d}{2} + k)\Gamma(\frac{k}{2})}(\tfrac{r}{2})^{k}\int_{0}^1 u^{k/2 - 1} {_0F_1}(\tfrac{d}{2} + k; -(\tfrac{r}{2})^2u) \, du.
\end{equation*}
Thus, \href{https://dlmf.nist.gov/16.5.E2}{equation~16.5.2} in~\cite{dlmf} applied with $a_0=k/2,$ $b_0=k/2+1$, $b_1=d/2+k$ and $z=-(\tfrac{r}{2})^2$ gives
\[
    \tm(r)=1-\frac{\Gamma(\frac{d + k}{2})}{\Gamma(\frac{d}{2} + k)\Gamma(\frac{k}{2}+1)}(\tfrac{r}{2})^k {_1F_2}(\tfrac{k}{2}; \tfrac{d}{2} + k, \tfrac{k}{2} + 1; -(\tfrac{r}{2})^2).
\]
Finally, coming back to $m_k(r)=\tm(2\pi r)$ we reach \eqref{eq: mk 2}. The proof of Proposition \ref{pro: mk} is thus completed.







\end{proof}

\begin{thebibliography}{99}
\bibitem[dlmf]{dlmf} \emph{NIST Digital Library of Mathematical Functions.} \url{https://dlmf.nist.gov/}, Release 1.2.4 of 2025-03-15. F.\,W.\,J. Olver, A.\,B. Olde Daalhuis, D.\,W. Lozier, B.\,I. Schneider, R.\,F. Boisvert, C.\,W. Clark, B.\,R. Miller, B.\,V. Saunders, H.\,S. Cohl, and M.\,A. McClain, (eds).
\bibitem[grafak]{grafakos} L. Grafakos, ``Classical Fourier Analysis'' Graduate Texts in Mathematics vol. 249, third edition, Springer Science+Business Media New York, 2014.
\bibitem[stein]{stein} E. M. Stein, ``Singular integrals and differentiability properties of functions'', Princeton University Press, Princeton, 1970. %
\end{thebibliography}
\end{document}
